% Options for packages loaded elsewhere
\PassOptionsToPackage{unicode}{hyperref}
\PassOptionsToPackage{hyphens}{url}
\documentclass[
  11pt,
]{article}
\usepackage{xcolor}
\usepackage[margin=25mm]{geometry}
\usepackage{amsmath,amssymb}
\setcounter{secnumdepth}{-\maxdimen} % remove section numbering
\usepackage{iftex}
\ifPDFTeX
  \usepackage[T1]{fontenc}
  \usepackage[utf8]{inputenc}
  \usepackage{textcomp} % provide euro and other symbols
\else % if luatex or xetex
  \usepackage{unicode-math} % this also loads fontspec
  \defaultfontfeatures{Scale=MatchLowercase}
  \defaultfontfeatures[\rmfamily]{Ligatures=TeX,Scale=1}
\fi
\usepackage{lmodern}
\ifPDFTeX\else
  % xetex/luatex font selection
\fi
% Use upquote if available, for straight quotes in verbatim environments
\IfFileExists{upquote.sty}{\usepackage{upquote}}{}
\IfFileExists{microtype.sty}{% use microtype if available
  \usepackage[]{microtype}
  \UseMicrotypeSet[protrusion]{basicmath} % disable protrusion for tt fonts
}{}
\makeatletter
\@ifundefined{KOMAClassName}{% if non-KOMA class
  \IfFileExists{parskip.sty}{%
    \usepackage{parskip}
  }{% else
    \setlength{\parindent}{0pt}
    \setlength{\parskip}{6pt plus 2pt minus 1pt}}
}{% if KOMA class
  \KOMAoptions{parskip=half}}
\makeatother
\setlength{\emergencystretch}{3em} % prevent overfull lines
\providecommand{\tightlist}{%
  \setlength{\itemsep}{0pt}\setlength{\parskip}{0pt}}
\usepackage{bookmark}
\IfFileExists{xurl.sty}{\usepackage{xurl}}{} % add URL line breaks if available
\urlstyle{same}
\hypersetup{
  hidelinks,
  pdfcreator={LaTeX via pandoc}}

\author{}
\date{}

\begin{document}

\section{Fredholm modules and analytic
K-homology}\label{fredholm-modules-and-analytic-k-homology}

\emph{Self-checked by the writing AI, GPT-6.1 Sol (OpenAI). Public
domain (CC0).}

The shift obstruction in extension theory can be recorded before taking
a quotient. A representation supplies the algebra action; an operator
that commutes with this action modulo compacts supplies the obstruction.
A Fredholm module is this pair of data, together with a parity. The
parity determines whether we will test it against projections or
unitaries.

Throughout, \(A\) is a separable complex C*-algebra and Hilbert spaces
are separable. Representations may be degenerate. Compact operators
between two Hilbert spaces mean norm limits of finite-rank operators. We
use the Fredholm criterion, index invariance under compact perturbations
and norm-continuous Fredholm paths, and index additivity proved in
Lemmas 1.1--1.2 of
\href{KT-KK-01.html\#lemma-1-1-the-fredholm-obstruction-survives-compact-errors}{the
extension lesson}. The preceding two lessons supply the Busby
correspondence and the complementary-corner construction for a
completely positive lift.

\subsection{1. The cycle and its three compact
defects}\label{the-cycle-and-its-three-compact-defects}

An \textbf{odd Fredholm module} over \(A\) is a triple \((H,\pi,F)\),
where \(\pi:A\to\mathcal B(H)\) is a *-representation and \(F\) is
bounded, such that, for every \(a\in A\),

\[
\begin{aligned}
\left[F,\pi(a)\right]&\in\mathcal K(H),\\
\pi(a)(F^2-1)&\in\mathcal K(H),\\
\pi(a)(F-F^*)&\in\mathcal K(H).
\end{aligned}
\tag{1.1}
\]

An \textbf{even Fredholm module} has in addition an orthogonal grading
\(H=H_+\oplus H_-\). Its grading operator \(\Gamma\) is \(1\) on \(H_+\)
and \(-1\) on \(H_-\). We require \(\Gamma\pi(a)=\pi(a)\Gamma\) and
\(\Gamma F=-F\Gamma\). Thus

\[
\begin{aligned}
\pi(a)&=\begin{pmatrix}\pi_+(a)&0\\0&\pi_-(a)\end{pmatrix},\\
F&=\begin{pmatrix}0&F_-\\F_+&0\end{pmatrix}.
\end{aligned}
\tag{1.2}
\]

The module is \textbf{degenerate} if all three expressions in (1.1) are
zero, for every \(a\). Such a module contributes no obstruction. In
particular, a self-adjoint unitary commuting with the representation is
degenerate. The zero representation makes every bounded \(F\)
degenerate.

For nonunital \(A\), the last two defects need not themselves be
compact. Their compactness is required only where the algebra acts. This
distinction is necessary, for example, when a cycle has an
infinite-dimensional zero-representation summand.

\textbf{Unitary equivalence} means a unitary intertwining both \(\pi\)
and \(F\); in the even case it must also preserve the grading. An
\textbf{operator homotopy} has fixed \(H,\pi\), and grading, and a
norm-continuous path \(F_t\) of operators satisfying (1.1). Direct sums
use block representations and operators, and the sum grading in the even
case.

We quotient cycles by the equivalence relation generated by unitary
equivalence, operator homotopy, and the relation identifying a cycle
with its direct sum with a degenerate cycle. We also identify every
degenerate cycle with the zero cycle. Write the resulting sets as
\(K^0(A)\) and \(K^1(A)\). Superscripts refer to analytic K-homology;
subscripts \(K_0(A),K_1(A)\) refer to the K-theory that will test it.

The comparison with the general Hilbert-module homotopy definition of
\(KK^i(A,\mathbb C)\) is a theorem, rather than part of the meaning of
operator homotopy. Its full proof must compare Hilbert-module homotopy
with the operator-homotopy and degenerate-summand relations used here.
That comparison remains unproved here.

\subsection{2. Normalization without a unitality
assumption}\label{normalization-without-a-unitality-assumption}

Before writing \(P=(1+F)/2\), we must make \(F\) a self-adjoint
involution. We develop the quotient-ideal and doubled-operator
construction using Blackadar, Sections 17.4--17.5. Its operators must be
checked against the locally compact defects in (1.1), since neither a
unit nor a nondegenerate representation is assumed here.

For a fixed representation put

\[
\begin{aligned}
\mathcal D_\pi&=\{T\in\mathcal B(H):[T,\pi(a)]\in\mathcal K(H)
                                  \text{ for all }a\},\\
\mathcal J_\pi&=\{T\in\mathcal D_\pi:T\pi(a)\in\mathcal K(H)
                                  \text{ for all }a\}.
\end{aligned}
\tag{2.1}
\]

\subsubsection{Lemma 2.1. The algebra of compact
defects}\label{lemma-2.1.-the-algebra-of-compact-defects}

\(\mathcal D_\pi\) is a unital C*-algebra and \(\mathcal J_\pi\) is a
closed two-sided *-ideal in it. For \(T\in\mathcal D_\pi\), membership
in \(\mathcal J_\pi\) is equivalent to compactness of \(\pi(a)T\) for
all \(a\). The cycle conditions say exactly that the image of \(F\) in
\(\mathcal D_\pi/\mathcal J_\pi\) is a self-adjoint involution.

\textbf{Proof.} For each \(a\), the commutator map
\(T\mapsto[T,\pi(a)]\) is norm-continuous and the compact ideal is norm
closed. The intersection of its inverse images is therefore closed. The
identity belongs to it, and the identities
\([ST,\pi(a)]=S[T,\pi(a)]+[S,\pi(a)]T\) and
\([T^*,\pi(a)]=-[T,\pi(a^*)]^*\) prove multiplication and adjoint
closure. It is consequently a unital C*-subalgebra of \(\mathcal B(H)\).

If \(T\in\mathcal D_\pi\), the equality \(\pi(a)T=T\pi(a)-[T,\pi(a)]\)
proves the left-right equivalence. The maps \(T\mapsto T\pi(a)\) are
also continuous, so their compact inverse images form a closed subspace.
If \(T\) lies in that subspace, \(T^*\pi(a)=(\pi(a^*)T)^*\) is compact.
For \(S\in\mathcal D_\pi\), \(ST\pi(a)\) is compact by the compact ideal
property, and \(TS\pi(a)=T\pi(a)S+T[S,\pi(a)]\) is compact as well.
These calculations prove exactly the closed two-sided *-ideal assertion.
The adjoint and square defects of a cycle belong to this ideal by the
left-right equivalence, so its quotient image is a self-adjoint element
whose square is the identity. Conversely those quotient equalities are
precisely the required locally compact defects. \(\square\)

\subsubsection{Proposition 2.2. Locally compact
perturbations}\label{proposition-2.2.-locally-compact-perturbations}

If \(F,F'\) are cycle operators for the same representation and parity,
and \((F'-F)\pi(a)\) is compact for every \(a\), then the segment
\(F_t=(1-t)F+tF'\) is an operator homotopy. In particular an ordinary
compact perturbation gives an operator homotopy.

\textbf{Proof.} Put \(K=F'-F\). Both endpoints lie in
\(\mathcal D_\pi\), and the assumed compact products put \(K\) in
\(\mathcal J_\pi\). For \(F_t=F+tK\), \[
F_t^2-1=(F^2-1)+t(FK+KF)+t^2K^2,
\qquad
F_t-F_t^*=(F-F^*)+t(K-K^*).
\] Every term on the right belongs to \(\mathcal J_\pi\), using the
ideal property for the products; their multiplication on the represented
algebra is therefore compact. The commutator is
\([F,\pi(a)]+t[K,\pi(a)]\), a sum of compacts. This proves every cycle
condition without imposing global compactness on \(K\). The path is
norm-continuous and is odd for the grading whenever its endpoints are.
An ordinary compact \(K\) satisfies the assumed products. \(\square\)

\subsubsection{Proposition 2.3. An exact involution
representative}\label{proposition-2.3.-an-exact-involution-representative}

Every even or odd module is equivalent to one with \(F=F^*\) and
\(F^2=1\). The construction can be performed continuously on a
norm-continuous family of cycle operators.

\textbf{Proof.} Lemma 2.1 permits functional calculus in the quotient
\(\mathcal D_\pi/\mathcal J_\pi\). Set \(F_s=(F+F^*)/2\). It has the
same quotient image as \(F\), so the segment joining them is allowed by
Proposition 2.2. That quotient image is an involution. For the odd
continuous clipping function \(g(t)=\max(-1,\min(t,1))\), put
\(f=g(F_s)\). Its quotient image is still the same involution, since
\(g(1)=1\) and \(g(-1)=-1\); hence \(f-F_s\in\mathcal J_\pi\) and a
second allowed segment joins them. The operator \(f\) is a self-adjoint
contraction. If the starting operator is odd, conjugating by the grading
replaces \(F_s\) by \(-F_s\), and the oddness of \(g\) makes \(f\) odd
too.

The positive operator \(s=(1-f^2)^{1/2}\) commutes with \(f\). Its
quotient image is zero, so \(s\in\mathcal J_\pi\). On \(H\oplus H\),
with representation \(\pi\oplus0\), form

\[
G_t=\begin{pmatrix}f&ts\\ts&-f\end{pmatrix},
\qquad 0\leq t\leq1.
\tag{2.2}
\]

For an even cycle use the grading
\(\operatorname{diag}(\Gamma,-\Gamma)\). The square root \(s\) is even,
because \(f^2\) is even, and the opposite second grading therefore makes
its off-diagonal positions odd. All \(G_t\) are self-adjoint. Direct
multiplication, using \(fs=sf\), gives \[
G_t^2-1=(t^2-1)
\begin{pmatrix}s^2&0\\0&s^2\end{pmatrix}.
\] Its product with the represented algebra is compact because
\(s^2\in\mathcal J_\pi\). The only other commutator blocks are
\([f,\pi(a)]\), \(-t\pi(a)s\) and \(ts\pi(a)\); these are compact by
\(f\in\mathcal D_\pi\) and \(s\in\mathcal J_\pi\). Thus (2.2) is an
operator homotopy. At \(t=0\) the first block is the clipped cycle and
the second is a zero-representation degenerate cycle. At \(t=1\) the
displayed square defect is exactly zero. These successive segments and
the degenerate-summand relation produce the required representative.

For a norm-continuous original family the self-adjoint parts are
uniformly bounded on compact parameter intervals. Polynomial
approximation to \(g\) on one common spectral interval proves norm
continuity of the clipping; polynomial approximation to the square root
on \([0,1]\) proves norm continuity of \(s\). Matrix assembly then
proves continuity of \(G_1\). This argument includes operators whose
spectra cross zero; no gap is used. \(\square\)

This construction deliberately allows a zero-representation summand. For
an even cycle of nonzero index, one cannot in general obtain an exact
odd involution by a compact perturbation on the original graded space:
an exact odd involution is an isomorphism between its two graded
summands.

\subsection{3. Addition, inverses and
pullback}\label{addition-inverses-and-pullback}

\subsubsection{Theorem 3.1. Analytic K-homology is an abelian
group}\label{theorem-3.1.-analytic-k-homology-is-an-abelian-group}

Direct sum makes \(K^0(A)\) and \(K^1(A)\) abelian groups. An inverse of
an odd cycle is \((H,\pi,-F)\). An inverse of an even cycle is
\((H^{\mathrm{op}},\pi,-F)\), where \(H^{\mathrm{op}}\) has the opposite
grading.

\textbf{Proof.} Direct sum preserves the defects and all generating
relations. Rebracketing and flipping summands are grading-preserving
unitaries where appropriate, giving associativity and commutativity. The
zero cycle is an identity by the degenerate-summand relation.

We may first replace \(F\) by its self-adjoint part, by Proposition 2.2.
Add the proposed inverse and use, on the two copies,

\[
F_\theta=
\begin{pmatrix}
\cos\theta\,F&\sin\theta\,1\\
\sin\theta\,1&-\cos\theta\,F
\end{pmatrix},
\tag{3.1}
\]

for \(0\leq\theta\leq\pi/2\), with representation \(\pi\oplus\pi\). In
the even case the opposite second grading makes the off-diagonal
identity odd. Squaring gives \[
F_\theta^2-1
=\cos^2\theta
\begin{pmatrix}F^2-1&0\\0&F^2-1\end{pmatrix}.
\] Its product with the representation is compact. The commutator has
diagonal blocks \(\cos\theta[F,\pi(a)]\) and their negatives and zero
off-diagonal blocks; self-adjointness is exact. At \(\theta=0\) we have
the cycle plus its proposed inverse. At \(\pi/2\) we have the
self-adjoint exchange unitary, commuting exactly with the
representation. It is degenerate, so the sum represents zero. Finally
\(-F\) and the grading reversal respect each generating equivalence
relation, so the inverse is well defined on classes. \(\square\)

\subsubsection{Proposition 3.2.
Contravariance}\label{proposition-3.2.-contravariance}

A *-homomorphism \(\alpha:A\to B\) induces \[
\begin{aligned}
\alpha^*&:K^i(B)\longrightarrow K^i(A),\\
[(H,\pi,F)]&\longmapsto[(H,\pi\alpha,F)].
\end{aligned}
\] It is a group homomorphism, and \((\beta\alpha)^*=\alpha^*\beta^*\),
\(1_A^*=1\).

\textbf{Proof.} Each compact condition over \(A\) is the corresponding
condition over \(B\), evaluated at \(\alpha(a)\). Degeneracy, unitary
equivalence and operator homotopy are preserved, as is direct sum.
Composition and the identity follow by composing the representations.
\(\square\)

When \(A\) is unital, one may also remove a degenerate part of its
representation. Put \(p=\pi(1)\). The off-diagonal blocks of \(F\)
relative to \(pH\oplus(1-p)H\) are compact, because \([F,p]\) is
compact. Delete them by Proposition 2.2. The remaining
zero-representation part is degenerate, and the part on \(pH\) has a
unital representation. The grading respects this decomposition for even
cycles. This reduction does not require a nondegeneracy convention at
the start.

\subsection{4. From a differential operator to a bounded
cycle}\label{from-a-differential-operator-to-a-bounded-cycle}

We state a spectral argument with all its operator hypotheses. It will
apply to Dirac operators and to compact-resolvent spectral triples.

\subsubsection{Lemma 4.0. Spectral calculus at compact
resolvent}\label{lemma-4.0.-spectral-calculus-at-compact-resolvent}

Let \(D\) be a densely defined self-adjoint operator with compact
\((1+D^2)^{-1}\). There is an orthonormal basis \(e_j\) of eigenvectors,
with real eigenvalues \(\lambda_j\), and \[
\operatorname{Dom}D=\left\{\sum_j u_je_j:\sum_j(1+\lambda_j^2)|u_j|^2<\infty\right\},
\qquad D\left(\sum_j u_je_j\right)=\sum_j\lambda_ju_je_j .
\] Only finitely many eigenvalues, counted with multiplicity, lie in any
bounded interval. Bounded scalar functions define bounded diagonal
operators \(f(D)\), preserving products and adjoints. Uniformly bounded
pointwise convergence of such functions on the eigenvalues gives strong
operator convergence. If \(f(t)\to0\) as \(|t|\to\infty\), then \(f(D)\)
is compact.

\textbf{Proof.} For \(\alpha>0\),
\(\|(D\pm i\alpha)u\|^2=\|Du\|^2+\alpha^2\|u\|^2\). Closedness of \(D\)
makes these ranges closed. Their orthogonal complements are the kernels
of \(D\mp i\alpha\), which are zero by the same equality and
self-adjointness. Thus both ranges are all \(H\), and their inverse
norms are at most \(1/\alpha\). Set \(R=(D-i)^{-1}\). Testing the
adjoint identity on the two domains gives \(R^*=(D+i)^{-1}\), and the
resolvent identity gives \[
R-R^*=2iRR^*=2iR^*R.
\] Hence \(T=RR^*\) is positive, injective, and has dense range. The
product maps into \(\operatorname{Dom}D^2\):
\(R^*v\in\operatorname{Dom}D\), \(DRw=w+iRw\), and \(R\) maps
\(\operatorname{Dom}D\) into \(\operatorname{Dom}D^2\). Multiplying out
on this domain proves \(T=(1+D^2)^{-1}\), so \(T\) is compact.

Here is the compact positive spectral argument used, including
completeness. If \(T\ne0\), put \(r=\|T\|\). The self-adjoint norm
identity in
\href{https://kokunoyumeto.github.io/open-math-courses/courses/foundations-of-von-neumann-algebras/hilbert-spaces-and-compact-operators.html}{Hilbert
spaces and compact operators}, Corollary 3.2, and positive functional
calculus give \(\sup_{\|u\|=1}\langle u,Tu\rangle=r\) and
\(0\leq T^2\leq rT\). Choose unit \(u_n\) with
\(\langle u_n,Tu_n\rangle\to r\). Then \[
\|(T-r)u_n\|^2
\leq r^2-r\langle u_n,Tu_n\rangle\longrightarrow0.
\] Compactness of \(T\) gives a convergent subsequence of \(Tu_n\),
hence of \(u_n\), whose unit limit is an eigenvector of eigenvalue
\(r\). Each nonzero eigenspace is finite dimensional: infinitely many
orthonormal vectors in one such space would have images under \(T\) with
no convergent subsequence. The orthogonal complement of any collected
eigenspaces reduces \(T\), and, if its restriction is nonzero, the same
argument supplies another positive eigenvalue. There can be only
finitely many eigenvectors with eigenvalues at least any fixed
\(\delta>0\), by compactness. Taking the closed span of all positive
eigenspaces leaves a reducing complement on which \(T=0\), because a
nonzero restriction would again supply a positive eigenvector.
Injectivity makes that complement zero.

Every positive eigenspace of \(T\) is contained in
\(\operatorname{Dom}D^2\), since \(v=\theta^{-1}Tv\). It is invariant
under \(D\), because the resolvents, hence \(T\), commute with \(D\) on
its domain. The restriction of \(D\) to this finite-dimensional space is
self-adjoint and diagonalizable. Its eigenvectors satisfy
\(Te_j=(1+\lambda_j^2)^{-1}e_j\). The preceding compactness property
proves the finiteness assertion for bounded \(\lambda_j\).

For \(u\in\operatorname{Dom}D\), self-adjointness gives
\(\langle e_j,Du\rangle=\lambda_j\langle e_j,u\rangle\), so the
indicated weighted sum is finite. Conversely its finiteness makes finite
eigenvector truncations converge together with their \(D\)-images;
closedness puts the limit in the domain and proves the formula. Bounded
diagonal multipliers have the stated products and adjoints by their
coordinates. Strong convergence follows by choosing a finite coordinate
truncation of a vector first, using the uniform bound for its tail, and
then passing to the pointwise limit on the finite part. Vanishing at
infinity makes the norm of the remaining diagonal tail tend to zero; the
finite truncations are finite rank. This proves compactness.
Finite-dimensional \(H\) satisfies the same conclusions with a finite
basis. \(\square\)

\subsubsection{Theorem 4.1. The bounded
transform}\label{theorem-4.1.-the-bounded-transform}

Let \(D\) be self-adjoint on \(H\), with \((1+D^2)^{-1}\) compact.
Suppose a dense *-subalgebra \(\mathcal A\subset A\) satisfies \[
\pi(a)\operatorname{Dom}D\subset\operatorname{Dom}D
\tag{4.1}
\] for every \(a\in\mathcal A\), and every commutator \([D,\pi(a)]\)
extends to a bounded operator. Then \[
F=D(1+D^2)^{-1/2}
\tag{4.2}
\] defines an odd Fredholm module. If \(H\) is graded, \(\pi\) is even,
the domain of \(D\) is grading invariant, and \(D\) is odd on that
domain, it defines an even module.

\textbf{Proof.} Lemma 4.0 defines (4.2) as the bounded diagonal scalar
function \(t/\sqrt{1+t^2}\) of \(D\). It gives \[
F=F^*,\quad \|F\|\leq1,\quad
F^2-1=-(1+D^2)^{-1}\in\mathcal K(H).
\tag{4.3}
\] We prove the commutator assertion rather than assuming that bounded
commutators with \(D\) automatically pass through functional calculus.

Set \(\alpha(s)=\sqrt{1+s^2}\) and \(R_\pm(s)=(D\pm i\alpha(s))^{-1}\).
These operators are compact. Indeed they are \((1+D^2)^{-1/2}\) times a
bounded continuous function of \(D\); the square root of a positive
compact operator is compact. The domain invariance in (4.1), applied
also to \(a^*\), justifies the resolvent identity \[
[R_\pm(s),\pi(a)]=-R_\pm(s)[D,\pi(a)]R_\pm(s).
\tag{4.4}
\] For example apply both sides to a vector, use that each resolvent
takes it into \(\operatorname{Dom}D\), and multiply by \(D\pm i\alpha\);
the identity follows there and extends boundedly.

The scalar identity \[
\frac{t}{\sqrt{1+t^2}}
=\frac1\pi\int_0^\infty
\left(\frac1{t-i\alpha(s)}+\frac1{t+i\alpha(s)}\right)\,ds
\tag{4.5}
\] holds because its truncation at \(R\) is
\(\frac{2t}{\pi\sqrt{1+t^2}}\arctan(R/\sqrt{1+t^2})\). These scalar
truncations have absolute value at most one and converge pointwise. On
every finite interval the resolvents are norm-continuous, so their
operator integrals exist; applying the integral to each eigenvector
gives exactly the displayed scalar truncation. Lemma 4.0 now makes these
integrals converge strongly to \(F\). No unbounded spectral theorem is
assumed in this passage.

Taking their commutators is better behaved: (4.4) gives a compact
integrand of norm at most \[
\frac{2\|[D,\pi(a)]\|}{\pi(1+s^2)}.
\] It is norm-continuous in \(s\), and this bound is integrable. Thus
its integral converges in norm to a compact operator. The strong limit
of the same commutators is \([F,\pi(a)]\), so that operator is compact
for \(a\in\mathcal A\). Since \(\|[F,\pi(a)]\|\leq2\|F\|\|a\|\), density
proves compactness for all \(a\in A\).

In the graded case \(\Gamma D\Gamma=-D\) with its domain retained.
Functional calculus therefore gives \(\Gamma F\Gamma=-F\), since the
scalar function in (4.2) is odd. This proves the required parity and
completes the proof. \(\square\)

Any continuous real normalizing function \(\chi\), with
\(\chi(t)\to\pm1\) as \(t\to\pm\infty\), gives the same class, provided
it is odd when a grading is present. The difference from
\(t/\sqrt{1+t^2}\) vanishes at infinity. A continuous function of \(D\)
vanishing at infinity is compact: truncate it to a bounded spectral
interval, whose spectral projection is finite rank because
\((1+D^2)^{-1}\) is compact, and estimate the remaining norm by its
supremum. Thus the two transforms differ by a compact operator, and
Proposition 2.2 applies.

Let \(M\) be a closed smooth manifold and \(E\) a smooth Hermitian
vector bundle of finite rank. Fix any smooth positive density on \(M\),
which defines \(L^2(M,E)\). The following proofs apply to this density
and require no orientation.

The elementary measure inputs are the actual earlier
\href{https://kokunoyumeto.github.io/open-math-courses/courses/foundations-of-von-neumann-algebras/compact-and-trace-class-operators-the-predual-of-b-h-and-the-operator-topologies.html\#oa-fnd-lt-17}{sigma-finite
product and real-line measure lemma}, Lemma 0.1(a)--(c): Fubini,
interval measure and one-dimensional \(C_c^\infty\)-density in \(L^2\).
Orthogonal projections, Hilbert-space forms, orthonormal expansions and
Parseval are proved in
\href{https://kokunoyumeto.github.io/open-math-courses/courses/foundations-of-von-neumann-algebras/hilbert-spaces-and-compact-operators.html}{Hilbert
spaces and compact operators}, Theorems 2.2, 3.1 and 4.1. The
finite-dimensional chain rule and change-of-variables formula are used
in smooth coordinates. No elliptic regularity, parametrix or
pseudodifferential boundedness result is an input.

\subsubsection{Lemma 4.1a. Finite-chart Sobolev density and
compactness}\label{lemma-4.1a.-finite-chart-sobolev-density-and-compactness}

The space \(H^1(M,E)\) is equivalently the completion of smooth sections
in a finite-chart first-derivative norm or the space of \(L^2\) sections
having locally square-integrable weak first derivatives in every smooth
frame. These definitions and their equivalent norms are independent of
the finite charts and cutoffs. Smooth sections are dense in both
\(H^1(M,E)\) and \(L^2(M,E)\). The inclusion
\(H^1(M,E)\hookrightarrow L^2(M,E)\) is a norm limit of finite-rank
operators. Multiplication by any smooth scalar function or smooth bundle
endomorphism is bounded on \(H^1\), and smooth functions are uniformly
dense in \(C(M)\).

\textbf{Proof.} We first give the finite Fourier argument, including
completeness. On an interval of length \(\ell\), rescale the circle
exponentials to the normalized functions
\(\ell^{-1/2}e^{2\pi ikx/\ell}\). Integration proves orthonormality. The
nonnegative circle polynomial \[
K_N(t)=\frac1N\left|\sum_{j=0}^{N-1}e^{ijt}\right|^2
\] has normalized integral one and satisfies
\(K_N(t)\leq [N\sin^2(\delta/2)]^{-1}\) for \(\delta\leq |t|\leq\pi\).
The integral is one because only the equal-index terms survive; the
inequality follows by summing the geometric progression. For a
continuous periodic function, convolution with \(K_N\) is a
trigonometric polynomial. Splitting its difference from that function
into \(|t|<\delta\) and its complement, uniform continuity controls the
first part and the displayed bound controls the second. Thus
trigonometric polynomials are uniformly dense in the continuous periodic
functions. The earlier interval-density lemma gives
\(C_c^\infty(0,\ell)\)-density in \(L^2(0,\ell)\); these functions
extend periodically and continuously by zero. Uniform approximation
therefore proves completeness of the one-dimensional Fourier family.

Products give a complete orthonormal family on any periodic cube
\(Q=(-\ell/2,\ell/2)^n\). Here is the induction proving completeness. If
\(g\in L^2(Q)\) is orthogonal to every product exponential, for each
\(k'\in\mathbb Z^{n-1}\) form its partial Fourier coefficient in the
first \(n-1\) coordinates. Fubini and Cauchy--Schwarz put that
coefficient in the \(L^2\) space of the last coordinate. All its
one-dimensional coefficients vanish, so it is zero almost everywhere.
Intersect the countably many full-measure sets for \(k'\), and also the
full-measure set on which the sections of \(g\) are in \(L^2\). Each
remaining section has all its \((n-1)\)-dimensional coefficients zero,
hence vanishes by induction. Fubini gives \(g=0\). This proves
completeness for every finite \(n\), also for vector-valued functions
component by component.

Write \[
e_k(x)=\ell^{-n/2}e^{i\kappa_k\cdot x},
\qquad \kappa_k=\frac{2\pi}{\ell}k,
\qquad \widehat v(k)=\int_Q \overline{e_k(x)}\,v(x)\,dx .
\] Parseval identifies the periodic first-derivative space with \[
\|v\|_{H^1(Q)}^2
=\sum_{k\in\mathbb Z^n}(1+|\kappa_k|^2)\|\widehat v(k)\|^2
=\|v\|_2^2+\sum_{j=1}^n\|\partial_jv\|_2^2 .
\] To verify the weak-derivative assertion rather than assume it,
testing a weak derivative against \(e_k\) gives coefficients
\(i\kappa_{k,j}\widehat v(k)\). Conversely, if the weighted sum is
finite, the finite Fourier sums converge in \(L^2\) together with their
derivatives. Integration by parts against a smooth periodic test
function passes to these limits and proves that the derivative limits
are the weak derivatives. This also proves completeness of this space,
since the weighted sequence space is complete. A function supported in
the interior of \(Q\), extended by zero and periodically, has the same
weak-derivative interpretation: in the testing identity insert a smooth
cutoff equal to one near its support, so no boundary term is present.

For the global construction choose finitely many bundle charts with
coordinate cubes \(Q_j\) having compact closures inside their charts,
and smooth real functions \(\chi_j\) supported in their interiors, with
\(\sum_j\chi_j^2=1\). These cutoffs can be constructed without an
analytic existence theorem: in each chart a translated and rescaled
version of \(\exp(-1/(1-|x|^2))\) on \(|x|<1\), extended by zero, is
positive on a smaller ball. Finitely many such smaller balls cover
\(M\); divide the cutoffs by the positive smooth square root of the sum
of their squares. Smoothness at the ball boundary follows because every
derivative is an exponential times a rational expression, and
\(t^{-a}e^{-1/t}\to0\) for every \(a\) as \(t\downarrow0\), as follows
from the exponential series. Smooth local frames may be made orthonormal
by the finite-dimensional Gram--Schmidt formulas. On the compact
coordinate cubes the density is \(\rho_j(x)\,dx\), with
\(0<c_j\leq\rho_j\leq C_j\).

Let \(S_ju\) be the coordinate vector of \(\chi_ju\), extended
periodically on \(Q_j\), and for smooth \(u\) define \[
\|u\|_{H^1}^2=\sum_j\|S_ju\|_{H^1(Q_j)}^2 .
\] This norm controls the global \(L^2\) norm, since the local \(L^2\)
norms are equivalent to their density-weighted norms and
\(\sum_j\chi_j^2=1\). Its completion embeds injectively in \(L^2\).
Indeed, if a Cauchy sequence converges to zero in \(L^2\), each
coordinate vector converges to zero there; its derivative limit is zero
by integration by parts against compactly supported tests. Thus its
\(H^1\) limit is zero.

We check all changes of coordinates and reconstructions used next. On a
compact overlap, a localized frame change and coordinate pullback has
the form \(v\mapsto C(x)v(y(x))\), with a smooth compactly supported
matrix factor \(C\). The ordinary chain rule gives \[
\partial_\ell(C\,v\circ y)
=(\partial_\ell C)v\circ y
+ C\sum_m(\partial_\ell y_m)(\partial_m v)\circ y .
\] The matrix coefficients and first derivatives are bounded there, as
are the Jacobian of the coordinate change and its reciprocal. The
change-of-variables formula consequently bounds the output \(H^1\) norm
by a constant times the input norm. For a periodic \(H^1\) vector the
same formula follows by approximating it by its finite Fourier sums.
Products with the cutoffs have support away from the chart boundary, so
the zero extensions are legitimate weak-derivative extensions. Define
\(R_jv\) to be the section with local vector \(\chi_jv\), extended by
zero. The preceding estimate proves that \(R=\sum_jR_j\) is bounded from
the finite product of the coordinate \(H^1\) spaces into the global
completion. It is also bounded on the corresponding \(L^2\) spaces. On
smooth sections \(RSu=u\), because \(\sum_j\chi_j^2=1\), and this
identity extends by \(L^2\) limits.

These facts prove the claimed weak-derivative description. A completion
element has its local weak derivatives by the preceding limit argument.
Conversely, a section whose localized coordinate vectors have weak
derivatives in \(L^2\) has \(Su\) in the product \(H^1\) space, so
\(u=RSu\) belongs to the completion. The same chain-rule estimate, using
\(RSu=u\), compares any two finite-chart norms on smooth sections and
hence on the completions. A covariant-derivative definition gives the
same space: in a frame \(\nabla_j=\partial_j+C_j(x)\), with smooth
bounded matrices on the compact supports, and the two first-derivative
norms therefore bound one another.

Let \(\Pi_{j,R}\) retain the finitely many modes with
\(|\kappa_k|\leq R\) in \(Q_j\). Then \(T_R=\sum_jR_j\Pi_{j,R}S_j\) is a
finite-rank map with smooth output. The Fourier formula gives \[
\|(1-\Pi_{j,R})v\|_2
\leq(1+R^2)^{-1/2}\|v\|_{H^1(Q_j)} .
\] Boundedness of the reconstruction proves
\(\|u-T_Ru\|_{L^2}\leq C(1+R^2)^{-1/2}\|u\|_{H^1}\). This is the
asserted norm approximation of the inclusion by finite-rank operators.
For a fixed \(v\in H^1(Q_j)\), the weighted Fourier tail tends to zero
as well; the \(H^1\)-bounded reconstruction therefore proves
\(T_Ru\to u\) in \(H^1\). For \(u\in L^2(M,E)\), ordinary Fourier-tail
convergence and \(L^2\)-bounded reconstruction prove \(T_Ru\to u\) in
\(L^2\). This proves both smooth-density assertions.

Multiplication by a smooth matrix is \(H^1\)-bounded by the displayed
product rule and the compact-support coefficient bounds. To check the
final uniform-density assertion explicitly, take a nonnegative smooth
integral-one function \(\eta\) supported in the unit ball, obtained by
normalizing the exponential cutoff above, and put
\(\eta_\varepsilon(x)=\varepsilon^{-n}\eta(x/\varepsilon)\). For
\(f\in C(M)\), the local function \(\chi_jf\), extended by zero, is
continuous with compact support. Its convolutions with
\(\eta_\varepsilon\) are smooth and converge uniformly to it: the error
is at most the supremum of \(|(\chi_jf)(x-z)-(\chi_jf)(x)|\) for
\(|z|\leq\varepsilon\). Reconstructing with \(\chi_j\) gives smooth
global functions tending uniformly to \(\sum_j\chi_j^2f=f\). On a
zero-dimensional component all these spaces are finite dimensional and
the assertions are immediate; compactness permits only finitely many
such chart components. \(\square\)

\subsubsection{Lemma 4.1b. The first-order elliptic
estimate}\label{lemma-4.1b.-the-first-order-elliptic-estimate}

Let \(D:C^\infty(M,E)\to C^\infty(M,E)\) be any smooth first-order
elliptic differential operator. There are constants \(C,C'>0\) such that
\[
\|u\|_{H^1}\leq C(\|Du\|_2+\|u\|_2),
\qquad \|Du\|_2\leq C'\|u\|_{H^1},
\qquad u\in C^\infty(M,E).
\] Neither estimate assumes a Dirac-type symbol or invertibility of
\(D\).

\textbf{Proof.} In a local frame write
\(P=\sum_{j=1}^n A_j(x)\partial_j+B(x)\), with smooth matrices. Our
principal-symbol convention is \(\sigma_P(x,\xi)=i\sum_jA_j(x)\xi_j\).
At a fixed point \(x_0\), ellipticity and compactness of the unit
spheres in the finite-dimensional covector and vector spaces give a
constant \(c>0\) such that \[
\left\|i\sum_jA_j(x_0)\xi_j v\right\|
\geq c|\xi|\|v\| .
\] Indeed the continuous left side on the product of the two unit
spheres is positive and attains a positive minimum. Let
\(P_0=\sum_jA_j(x_0)\partial_j\). For a smooth vector \(v\) compactly
supported in a coordinate cube, extend it periodically and use Lemma
4.1a. The Fourier coefficient of \(P_0v\) is
\(i\sum_jA_j(x_0)\kappa_{k,j}\widehat v(k)\). Squaring the preceding
inequality and summing by Parseval gives
\(\|\nabla v\|_2\leq c^{-1}\|P_0v\|_2\).

Shrink the coordinate neighborhood \(W\) of \(x_0\) until
\((\sum_j\|A_j(x)-A_j(x_0)\|^2)^{1/2}\leq c/2\) on it. For
\(v\in C_c^\infty(W)\), Cauchy--Schwarz in the finite sum gives \[
\begin{aligned}
\|\nabla v\|_2
&\leq c^{-1}\big(\|Pv\|_2
+(c/2)\|\nabla v\|_2+\|B\|_{\infty,W}\|v\|_2\big),\\
\|\nabla v\|_2
&\leq 2c^{-1}\big(\|Pv\|_2+\|B\|_{\infty,W}\|v\|_2\big).
\end{aligned}
\] This absorbs the coefficient variation and proves the local estimate.
Density weights and smooth-frame norms change its constant only by their
bounded upper and lower factors.

Choose finitely many such neighborhoods covering \(M\), with cutoffs
supported compactly inside them as in Lemma 4.1a. The norm defined from
this refined cover is equivalent to the earlier \(H^1\) norm. Apply the
local estimate to each \(\chi_ju\), using \[
D(\chi_ju)=\chi_jDu+[D,\chi_j]u,
\qquad [D,\chi_j]=\sum_\ell A_\ell\,\partial_\ell\chi_j .
\] The commutator is a smooth bounded multiplication operator. Summing
the finitely many estimates gives the first global inequality. In the
same charts the bounded coefficients and the product rule give the
second inequality, by reconstructing \(u=\sum_j\chi_j(\chi_ju)\). A
zero-dimensional component contributes only a finite-dimensional norm
bound and the harmless \(\|u\|_2\) term. This proves both estimates in
the full stated generality. \(\square\)

\subsubsection{Lemma 4.1c. Friedrichs regularization and the weak
domain}\label{lemma-4.1c.-friedrichs-regularization-and-the-weak-domain}

For the operator \(D\) in Lemma 4.1b, \[
u\in L^2(M,E),\quad Du\in L^2(M,E)
\text{ in the sense of distributions}
\quad\Longleftrightarrow\quad u\in H^1(M,E).
\] The forward implication has the same bound as Lemma 4.1b. Locally the
mollifier commutator extends to uniformly bounded operators on \(L^2\)
and tends strongly to zero. In particular smooth approximations can be
chosen to converge in the graph norm.

\textbf{Proof.} We supply the commutator bound and convergence before
using them. On a fixed coordinate neighborhood extend the local
coefficients smoothly to \(\mathbb R^n\), agreeing on a larger
neighborhood of the compact support under consideration and having
bounded coefficients and first derivatives; multiply by a smooth cutoff
inside the original chart to do this. These extensions need not be
elliptic outside that neighborhood. Let \[
P=\sum_jA_j(x)\partial_j+B(x),\qquad
J_\varepsilon u=\eta_\varepsilon*u ,
\] where \(\eta\) is the nonnegative smooth integral-one ball kernel
constructed in Lemma 4.1a. For a nonnegative integrable kernel \(k\) of
mass \(m\), Cauchy--Schwarz in its measure gives \[
\left|\int k(z)v(x-z)\,dz\right|^2
\leq m\int k(z)|v(x-z)|^2\,dz .
\] Integration and Fubini give \(\|k*v\|_2\leq m\|v\|_2\). The same
estimate applied to norms controls a matrix kernel bounded in norm by
\(k\). Thus \(\|J_\varepsilon\|\leq1\). For a compactly supported vector
in the interior of a fixed cube and small \(\varepsilon\), convolution
equals periodic convolution there; its Fourier multiplier is
\(\int\eta(z)e^{-i\varepsilon\kappa_k\cdot z}\,dz\), bounded by one and
tending to one for each \(k\). Finite Fourier truncation followed by a
tail estimate proves \(J_\varepsilon u\to u\) in \(L^2\), and also in
\(H^1\) when \(u\in H^1\). It gives a smooth function: each derivative
falls on the smooth kernel, and Cauchy--Schwarz on its compact support
and dominated differentiation justify the integral derivatives.
Compactly supported smooth vectors are dense in \(L^2(\mathbb R^n)\):
first truncate a vector to an inner cube, use its Fourier approximation
in a larger cube, and multiply the finite Fourier sums by a smooth
cutoff equal to one on the inner cube. The cutoff does not increase the
\(L^2\) error. Consequently the contraction bound extends the \(L^2\)
convergence to every vector.

For a compactly supported smooth \(u\), integration by parts in the
convolution gives the exact formula \[
\begin{aligned}
R_\varepsilon u(x)
&:=(PJ_\varepsilon-J_\varepsilon P)u(x)\\
&=\sum_j\int
\big(A_j(x)-A_j(y)\big)
(\partial_j\eta_\varepsilon)(x-y)u(y)\,dy\\
&\quad+\sum_j\int
\eta_\varepsilon(x-y)(\partial_jA_j)(y)u(y)\,dy\\
&\quad+\int
\big(B(x)-B(y)\big)\eta_\varepsilon(x-y)u(y)\,dy .
\end{aligned}
\] The sign of the second line follows from
\(\partial_{y_j}\eta_\varepsilon(x-y)=-\partial_{x_j}\eta_\varepsilon(x-y)\);
integration by parts differentiates \(A_j(y)\) as well. Put
\(L_j=\sup_x\|dA_j(x)\|\) and \(M_j=\|\partial_jA_j\|_\infty\), where
the first norm is for the Euclidean differential into the matrix norm.
The mean-value formula bounds the first integrand's matrix kernel by
\(L_j|x-y||\partial_j\eta_\varepsilon(x-y)|\). Its scalar kernel has
mass \(L_j\int |z||\partial_j\eta(z)|\,dz\), independent of
\(\varepsilon\). The second kernel has mass at most \(M_j\), and the
last at most \(2\|B\|_\infty\). The preceding convolution estimate
therefore proves \[
\|R_\varepsilon u\|_2\leq C_P\|u\|_2,\qquad
C_P=\sum_j\left(
L_j\int |z||\partial_j\eta(z)|\,dz+M_j\right)
+2\|B\|_\infty .
\] Every integral in this constant is finite. The kernel formula defines
a bounded operator on all \(L^2\), with this same constant. The identity
with the distributional commutator holds there as well: approximate by
compactly supported smooth vectors in \(L^2\); multiplication and
differentiation converge on every compactly supported distributional
test, and the convolution kernel has compact support. Thus the smooth
identity passes to the distributional limit.

If \(v\) is compactly supported and smooth, \(J_\varepsilon v\to v\) in
\(H^1\), the bounded-coefficient operator
\(P:H^1(\mathbb R^n)\to L^2(\mathbb R^n)\) is bounded by its explicit
first-derivative formula, and \(J_\varepsilon Pv\to Pv\) in \(L^2\).
Hence \(R_\varepsilon v\to0\). For arbitrary \(u\in L^2\), approximate
it by such a \(v\) and use \[
\|R_\varepsilon u\|_2
\leq C_P\|u-v\|_2+\|R_\varepsilon v\|_2 .
\] Letting first \(\varepsilon\to0\) and then \(v\to u\) proves strong
convergence to zero. This proves the full uniformly bounded Friedrichs
commutator assertion, including its distributional interpretation.

Now let \(u\in L^2(M,E)\) and \(Du=f\in L^2\) distributionally. Use the
refined cutoffs from Lemma 4.1b. In one chart the compactly supported
vector \(v=\chi_ju\) satisfies \[
Pv=\chi_jf+[D,\chi_j]u\in L^2 .
\] The product rule here follows directly by testing
\(\partial(\chi_ju)=\chi_j\partial u+(\partial\chi_j)u\) against a
compactly supported smooth test function. The equation is a local
coordinate distributional equation. Indeed in an orthonormal frame with
density \(\rho\,dx\) the formal adjoint is
\(P^\dagger\phi=-\rho^{-1}\sum_j\partial_j(\rho A_j^*\phi)+B^*\phi\);
ordinary integration by parts proves that testing against this adjoint
is exactly the coordinate distributional equation, using the invertible
smooth density factor. Extension by zero creates no boundary
distribution, because \(\operatorname{supp}\chi_j\) is compactly inside
the chart.

Choose the coefficient extensions above equal to the original ones on
the local estimate neighborhood \(W\). For sufficiently small
\(\varepsilon\), \(v_\varepsilon=J_\varepsilon v\) is smooth and
compactly supported in \(W\), and the proved commutator identity gives
\[
Pv_\varepsilon=J_\varepsilon(Pv)+R_\varepsilon v
\longrightarrow Pv\quad\text{in }L^2,
\qquad v_\varepsilon\longrightarrow v\quad\text{in }L^2 .
\] The local elliptic estimate of Lemma 4.1b, applied to
\(v_\varepsilon-v_\delta\), makes this family Cauchy in \(H^1\) as
\(\varepsilon,\delta\downarrow0\). Completeness and the \(L^2\)
identification in Lemma 4.1a show that its \(H^1\) limit is \(v\). Thus
each \(\chi_ju\) has square-integrable weak derivatives, and
reconstruction puts \(u\) in \(H^1(M,E)\). Passing to the limit in the
local estimates and summing gives the global bound of Lemma 4.1b for
this \(u\). Conversely, smooth density in \(H^1\) and the bounded map
\(D:H^1\to L^2\) show that every \(H^1\) section has its distributional
\(D\)-image in \(L^2\). The same density gives convergence in the graph
norm. \(\square\)

\subsubsection{Corollary 4.2. Closed-manifold differential
operators}\label{corollary-4.2.-closed-manifold-differential-operators}

Let \(M\) be a closed smooth manifold and \(E\) a Hermitian vector
bundle. If a first-order elliptic differential operator \(D\) is
formally self-adjoint on smooth sections, its closure on \(L^2(M,E)\)
has domain \(H^1(M,E)\), is self-adjoint, and its bounded transform
defines a module over \(C(M)\). A bundle grading for which \(D\) is odd
makes the module even.

\textbf{Proof.} Smooth sections are \(L^2\)-dense by Lemma 4.1a. Formal
self-adjointness makes the smooth-domain operator closable: if
\(u_j\to0\) and \(Du_j\to g\) in \(L^2\), testing against every smooth
\(\phi\) gives
\(\langle g,\phi\rangle=\lim_j\langle u_j,D\phi\rangle=0\); density
gives \(g=0\). Its graph closure has domain exactly \(H^1\). In one
direction, smooth \(H^1\)-approximation and boundedness of
\(D:H^1\to L^2\) put \(H^1\) in that closure. In the other direction, a
graph-Cauchy smooth sequence is \(H^1\)-Cauchy by Lemma 4.1b, and its
\(L^2\) limit is the \(H^1\) limit from Lemma 4.1a.

The adjoint domain is exactly the set of \(L^2\) sections \(v\) for
which \(Dv\in L^2\) distributionally: testing the adjoint identity on
smooth sections and formal self-adjointness give this equivalence. Lemma
4.1c puts each such \(v\) in \(H^1\); conversely smooth
\(H^1\)-approximation passes the formal adjoint identity to every
\(v\in H^1\). Thus the adjoint and closure have the same domain and
action. This proves self-adjointness with no separate regularity
assumption.

For \(u\in\operatorname{Dom}D\), self-adjointness gives
\(\|(D-i)u\|^2=\|Du\|^2+\|u\|^2\). The range is closed by closedness of
\(D\), and its orthogonal complement is \(\ker(D+i)=0\). Hence
\(Q=(D-i)^{-1}\) exists, has norm at most one on \(L^2\), and maps
\(L^2\) boundedly into \(H^1\) by Lemma 4.1b. Put \(K=-iQ\). Since
\(QDu=u+iQu\) for \(u\in H^1\), we retain the precise identity and
estimate \[
u=QDu+Ku,\qquad
\|u\|_{H^1}\leq C(\|Du\|_{L^2}+\|u\|_{L^2}).
\tag{4.6}
\] Here \(Q,K:L^2\to H^1\) are the explicitly constructed resolvent
maps; this identity is used on \(H^1\), after the weak-domain proof, and
makes no claim about a pseudodifferential order. Compactness of
\(H^1\hookrightarrow L^2\), proved in Lemma 4.1a by finite Fourier
truncations, makes \(Q\) compact as an operator on \(L^2\). The
resolvent identity and adjoint give \(QQ^*=(1+D^2)^{-1}\), as in Lemma
4.0, so that operator is compact too.

Multiplication by a smooth function preserves \(H^1\) by Lemma 4.1a. In
local coordinates \(D=\sum_jA_j(x)\partial_j+B(x)\), so
\([D,f]=\sum_jA_j(x)(\partial_jf)\) is a smooth bundle multiplication
operator, bounded on the compact manifold. Lemma 4.1a also proves
uniform smooth density in \(C(M)\). Theorem 4.1 therefore applies with
every original hypothesis checked. A smooth bundle grading is
\(H^1\)-bounded by Lemma 4.1a, and its odd-domain identity extends from
smooth sections by density, so the graded conclusion follows as well.
\(\square\)

Formal self-adjointness is a real hypothesis here. Ellipticity by itself
gives a Fredholm realization between suitable Sobolev spaces, but does
not make a differential operator into the self-adjoint \(D\) required in
Theorem 4.1.

A compact-resolvent \textbf{spectral triple} consists of a dense
*-algebra \(\mathcal A\subset A\), a representation on \(H\), and a
self-adjoint \(D\) satisfying (4.1), such that every \([D,\pi(a)]\),
\(a\in\mathcal A\), extends to a bounded operator and \((1+D^2)^{-1}\)
is compact. These hypotheses give its bounded Fredholm module by Theorem
4.1. Summability, regularity and meromorphic dimension spectrum are
additional structures, studied in \emph{Spectral triples and dimension
spectrum}; none is needed for this construction. Locally compact triples
on noncompact spaces need the localized version of the bounded-transform
theorem, whose full proof remains required in \emph{Unbounded Kasparov
modules and spectral triples}; the compact-resolvent proof here does not
establish it.

\subsection{5. Three explicit examples}\label{three-explicit-examples}

\subsubsection{Lemma 5.0. The Fourier basis used in the
examples}\label{lemma-5.0.-the-fourier-basis-used-in-the-examples}

The functions \(e^{int}\), \(n\in\mathbb Z\), are an orthonormal basis
of the normalized \(L^2(S^1)\), and trigonometric polynomials are
uniformly dense in \(C(S^1)\). Products of these functions form an
orthonormal basis on the two-dimensional torus.

\textbf{Proof.} Integration of exponentials gives orthonormality. For
\(N\geq1\) use the nonnegative trigonometric polynomial \[
K_N(t)=\frac1N\left|\sum_{j=0}^{N-1}e^{ijt}\right|^2 .
\] Its integral with normalized angular measure is \(1\), since only the
equal-index terms survive. For \(\delta\leq|t|\leq\pi\), the
geometric-sum formula gives \(K_N(t)\leq1/(N\sin^2(\delta/2))\). For
continuous periodic \(f\), its convolution with \(K_N\) is a
trigonometric polynomial. Split \[
(K_N*f)(x)-f(x)=\int K_N(t)(f(x-t)-f(x))\,\frac{dt}{2\pi}
\] at \(|t|=\delta\). Uniform continuity makes the first part uniformly
small for small \(\delta\); the displayed tail bound makes the second
part uniformly small as \(N\to\infty\). This proves uniform density.
Using the product of the two kernels and splitting where either
coordinate is outside its small interval proves the same uniform
approximation for continuous functions on the two-dimensional torus.

The existing
\href{https://kokunoyumeto.github.io/open-math-courses/courses/foundations-of-von-neumann-algebras/compact-and-trace-class-operators-the-predual-of-b-h-and-the-operator-topologies.html\#oa-fnd-lt-17}{real-line
measure and density lemma}, Lemma 0.1(b)--(c), proves that endpoints
have measure zero and \(C_c^\infty(0,2\pi)\) is dense in
\(L^2(0,2\pi)\). Such functions extend continuously and periodically by
zero. Uniform trigonometric approximation therefore gives
\(L^2\)-density as well, proving completeness of the one-dimensional
orthonormal family.

For the two-dimensional claim, let \(g\in L^2((S^1)^2)\) be orthogonal
to all products. Fubini, proved in that same earlier Lemma 0.1(a), and
Cauchy--Schwarz show that \(g_n(y)=\int g(x,y)e^{-inx}\,dx/(2\pi)\)
belongs to \(L^2(S^1)\). All its Fourier coefficients vanish, so
\(g_n=0\) almost everywhere by the one-dimensional result. Intersect the
countably many full-measure sets for \(n\in\mathbb Z\). On their
intersection the \(L^2\) section \(g(\cdot,y)\) has every coefficient
zero, hence vanishes. Fubini gives \(g=0\) almost everywhere. This
proves the product family's completeness; orthonormality follows by
product integration. Parseval and graph-coordinate convergence are
proved in
\href{https://kokunoyumeto.github.io/open-math-courses/courses/foundations-of-von-neumann-algebras/hilbert-spaces-and-compact-operators.html}{Hilbert
spaces and compact operators}, Theorem 4.1. \(\square\)

\subsubsection{The Hardy module of the
circle}\label{the-hardy-module-of-the-circle}

Put \(H=L^2(S^1)\), with normalized angular measure, and let
\(e_n(z)=z^n\), \(n\in\mathbb Z\). Let \(\pi(f)\) be multiplication by
\(f\), and let \(P\) project onto the closed span of \(e_n\) for
\(n\geq0\). Then \[
F=2P-1
\tag{5.1}
\] is a self-adjoint involution. For \(k>0\), the commutator
\([P,\pi(z^k)]\) vanishes on every basis vector except
\(e_{-k},\ldots,e_{-1}\); its image lies in
\(\operatorname{span}(e_0,\ldots,e_{k-1})\). For \(k<0\) the same
conclusion follows by adjoints. Thus these commutators have finite rank.
Trigonometric polynomials approximate every continuous function
uniformly, and \(\|[P,\pi(f)]\|\leq2\|f\|_\infty\), so \([P,\pi(f)]\) is
compact for every \(f\in C(S^1)\). This verifies all module conditions.

The operator \(D=-i\,d/d\theta\) on the circle has \(De_n=ne_n\), with
domain \(\{\sum c_ne_n:\sum n^2|c_n|^2<\infty\}\). Testing an adjoint
vector against every Fourier basis vector gives the same square-summable
weighted domain, so this diagonal realization is self-adjoint; finite
Fourier sums converge in its graph norm and prove that it is the closure
of the smooth differential operator. Its bounded transform has
eigenvalues \(n/\sqrt{1+n^2}\). The differences from (5.1) tend to zero
as \(|n|\to\infty\); the discrepancy at \(n=0\) affects only one vector.
Their difference is compact. Hence the Hardy cycle represents the same
K-homology class as this bounded transform.

Compressing \(\pi(z)\) to \(PH\) gives the unilateral shift
\(Se_n=e_{n+1}\), \(n\geq0\). Its kernel is zero and its cokernel is
spanned by \(e_0\), so \[
\operatorname{index}(P\pi(z)P)=-1.
\tag{5.2}
\] We will prove below that this index detects a nonzero class and all
its multiples.

\subsubsection{A Dirac operator on the two-dimensional
torus}\label{a-dirac-operator-on-the-two-dimensional-torus}

Let \(T^2=(\mathbb R/2\pi\mathbb Z)^2\), take
\(H=L^2(T^2,\mathbb C^2)\), and grade by \[
\Gamma=\begin{pmatrix}1&0\\0&-1\end{pmatrix}.
\] Use the Pauli matrices \[
\sigma_1=\begin{pmatrix}0&1\\1&0\end{pmatrix},\qquad
\sigma_2=\begin{pmatrix}0&-i\\i&0\end{pmatrix},
\qquad D=-i(\sigma_1\partial_x+\sigma_2\partial_y).
\tag{5.3}
\] They are self-adjoint, square to one, and anticommute. On the Fourier
mode \(e^{i(nx+my)}v\), \(D\) acts by the self-adjoint matrix
\(n\sigma_1+m\sigma_2\), whose square is \((n^2+m^2)1\). Consequently
the closure is self-adjoint with domain \[
\left\{(v_{n,m}):\sum_{n,m}(1+n^2+m^2)\|v_{n,m}\|^2<\infty\right\}.
\] This can be verified directly: the direct sum of the
finite-dimensional self-adjoint matrices has adjoint with that same
domain, and finite Fourier sums are dense in the graph norm.

The eigenvalues of \((1+D^2)^{-1}\) are \((1+n^2+m^2)^{-1}\), of
multiplicity two, and tend to zero outside finite sets. Thus this
operator is compact. Multiplication by \(C(T^2)\) is even, \(D\)
anticommutes with \(\Gamma\), and smooth commutators are
\(-i(\sigma_1\partial_x f+\sigma_2\partial_y f)\). For domain
preservation, finite Fourier sums have graph convergence, and Parseval
identifies this domain with the functions whose two weak first
derivatives are in \(L^2\). For smooth \(f\),
\(\partial_j(fu)=f\partial_ju+(\partial_jf)u\); the bounded coefficients
give a graph-norm bound, first on those finite sums and then by
completion. Thus multiplication by \(f\) preserves the domain. Smooth
functions are uniformly dense by Lemma 5.0. All hypotheses of Theorem
4.1 have now been checked, and it gives an even module.

Its positive-to-negative bounded block on the \((n,m)\)-mode is \[
F_+(n,m)=\frac{n+im}{\sqrt{1+n^2+m^2}}.
\tag{5.4}
\] Both \(F_+\) and its adjoint have the one-dimensional constant-mode
kernel. On the other modes the modulus is at least \(1/\sqrt2\), so the
range is closed. Its untwisted index is therefore \(1-1=0\). This
computation concerns the test by the trivial line bundle; it does not
assert that the torus Dirac class vanishes.

\subsubsection{Finite-dimensional
cycles}\label{finite-dimensional-cycles}

If \(H\) is finite dimensional, every defect is compact, so every
representation and bounded operator of the appropriate parity is a
cycle. Such a cycle need not be degenerate: \(F=0\) and a nonzero
representation give \(\pi(a)(F^2-1)=-\pi(a)\).

Over \(A=\mathbb C\), an even finite-dimensional example is
\(H_+=\mathbb C^r\), \(H_-=\mathbb C^s\), scalar representation and
\(F=0\). Its eventual even index test is \(r-s\). A pair with \(r=s=1\)
is zero already here: the path \[
\begin{pmatrix}0&t\\t&0\end{pmatrix},
\qquad0\leq t\leq1,
\] has only finite-dimensional defects and ends at a commuting
self-adjoint exchange unitary. In contrast an odd finite-dimensional
cycle is always zero. Replace \(F\) along a segment by \(1\); all
defects remain compact, and the endpoint commutes with the
representation and is a self-adjoint unitary. This parity difference
explains why a single finite-dimensional graded representation can
encode an even index, whereas odd index information needs an
infinite-dimensional space.

\subsection{6. Odd cycles and semisplit
extensions}\label{odd-cycles-and-semisplit-extensions}

\subsubsection{Proposition 6.1. The compression
extension}\label{proposition-6.1.-the-compression-extension}

Let \((H,\pi,F)\) be an odd module with \(F=F^*\), \(F^2=1\), and put
\(P=(1+F)/2\). The formula \[
\tau(a)=q_{PH}\big(P\pi(a)|_{PH}\big)
\tag{6.1}
\] defines a Busby map into \(\mathcal Q(PH)\) with a completely
positive contractive lift. Its pullback is a semisplit extension of
\(A\) by \(\mathcal K(PH)\).

\textbf{Proof.} The inclusion \(V:PH\hookrightarrow H\) is an isometry
and the proposed lift is \(L(a)=V^*\pi(a)V\). For a positive matrix
\([a_{ij}]\), its representation on \(H^n\) is positive, and compression
by \(V^{(n)}\) gives \([L(a_{ij})]\geq0\); hence \(L\) is completely
positive. Contractivity of \(\pi\) and \(V\) proves contractivity of
\(L\), and taking adjoints gives \(L(a^*)=L(a)^*\). Its multiplication
error is exactly \[
L(ab)-L(a)L(b)
=P\pi(a)(1-P)\pi(b)|_{PH}.
\tag{6.2}
\] On \(PH\), \((1-P)\pi(b)P=(1-P)[\pi(b),P]\), a compact operator
between the two spectral ranges. Thus this error vanishes in the Calkin
quotient and \(q_{PH}L\) preserves multiplication as well as adjoints.
For its pullback algebra, \(a\mapsto(L(a),a)\) lies in the pullback, is
a completely positive contraction at every matrix level, and is a
section of the quotient. The Busby pullback theorem proved in the first
extension lesson therefore gives precisely a semisplit extension by
\(\mathcal K(PH)\). \(\square\)

This extension is essential precisely when the Busby map is injective;
the construction makes no automatic essentiality claim. If \(PH\) is
finite dimensional, the quotient is zero. To use a fixed
infinite-dimensional compact ideal in all cases, append an
infinite-dimensional zero-representation summand on which \(F=1\). It is
degenerate, and enlarges \(PH\) to \(PH\oplus\ell^2\); the new Busby map
is the old one plus the split zero map.

\subsubsection{Proposition 6.2. A semisplit extension gives an odd
cycle}\label{proposition-6.2.-a-semisplit-extension-gives-an-odd-cycle}

Every semisplit extension of separable \(A\) by \(\mathcal K(H)\) is
obtained by the compression construction from an odd module with an
exact self-adjoint involution.

\textbf{Proof.} Write \(\tau\) for the given Busby map. If \(H\) is
finite dimensional, \(\mathcal Q(H)=0\) and \(\tau=0\). Its zero lift is
the compression of the zero representation on \(H\oplus H\), with
\(P=\operatorname{diag}(1,0)\) and \(F=2P-1\). This is a degenerate
cycle. The first lesson's Busby correspondence identifies its pullback
with the given extension, so this case is included rather than excluded
by a fixed infinite-dimensional convention.

For infinite-dimensional separable \(H\), semisplitting provides a
completely positive contractive lift \(L:A\to\mathcal B(H)\). The
previous lesson, Lemma 2.2, proves the corner dilation from the strict
KSGNS construction and operator-range stabilization. For \(B=\mathbb C\)
it produces a homomorphism \[
\rho:A\longrightarrow\mathcal B(H\oplus H),\qquad
\rho(a)=\begin{pmatrix}L(a)&\rho_{12}(a)\\
                                  \rho_{21}(a)&\rho_{22}(a)\end{pmatrix}.
\] Multiplication of its first diagonal corner at \(a^*a\) gives \[
\rho_{21}(a)^*\rho_{21}(a)=L(a^*a)-L(a)^*L(a).
\] The quotient of the right side is zero since \(qL=\tau\) is
multiplicative. The C*-identity in the quotient forces
\(q\rho_{21}(a)=0\), hence compactness. Applying this to \(a^*\) and
using \(\rho_{12}(a)=\rho_{21}(a^*)^*\) proves compactness of the other
off-diagonal corner.

Now \(P=\operatorname{diag}(1,0)\) and \(F=2P-1\) have zero adjoint and
square defects, and \[
[F,\rho(a)]=\begin{pmatrix}0&2\rho_{12}(a)\\ -2\rho_{21}(a)&0\end{pmatrix}
\] is compact. This is the required exact odd cycle. Compressing its
representation to \(PH=H\) gives \(L\), hence its compression Busby map
is exactly \(\tau\). The Busby correspondence once more identifies the
extensions. No essentiality, nuclearity or unit was used; separability
is the hypothesis needed by that corner-stabilization step. \(\square\)

For the Hardy module, (6.1) is precisely the Toeplitz Busby map. The
same construction also explains the inverse extension: switch to the
negative spectral subspace, replacing \(F\) by \(-F\). The two
compressed corners together are the quotient of the full representation,
a split extension.

The full class-level statement is \[
K^1(A)\ \cong\ KK^1(A,\mathbb C)\
\cong\ \operatorname{Ext}(A)^{-1},
\tag{6.3}
\] where the last group is the subgroup of invertible elements in the
ordinary stabilized extension monoid. The correspondence sends an exact
odd involution to (6.1), and a semisplit extension to Proposition 6.2.
It is natural under pullback in \(A\), and the compression index is its
pairing with \(K_1(A)\).

Formula (6.3) retains its full separable-algebra statement, but is not
established by the two constructions above. A complete proof still has
to show that compression respects all class relations, that changes of
dilation and normalization give the same class, that the constructions
are inverse on classes, and that Hilbert-module homotopy gives exactly
the analytic equivalence relation. The pending lesson \emph{Pictures of
KK: Fredholm picture, quasihomomorphisms, \(KK^1\) and extensions} is
not a delivered proof provider. The constructions here use no planar
classification theorem. For nuclear \(A\), the preceding lesson
independently proves that every extension is semisplit, so every
extension has the cycle representative constructed in Proposition 6.2;
this alone does not prove (6.3).

\subsection{7. An index test that already detects the
circle}\label{an-index-test-that-already-detects-the-circle}

Here we need only a fixed unitary in a unital algebra. The next lesson
proves the full matrix-stable, bilinear pairing.

\subsubsection{Lemma 7.1. A unitary tests an odd K-homology
class}\label{lemma-7.1.-a-unitary-tests-an-odd-k-homology-class}

For unital \(A\) and \(u\in A\) unitary, the positive-compression index
defines a group homomorphism \[
I_u:K^1(A)\longrightarrow\mathbb Z.
\] For an exact involution and a unital representation it is
\(\operatorname{index}(P\pi(u)P|_{PH})\).

\textbf{Proof.} For a general cycle use the particular continuous
normalization \(N(F)=G_1\) in Proposition 2.3, with representation
\(\rho=\pi\oplus0\), and put \(P_F=(1+N(F))/2\). The operator \[
U=\rho(u)+(1-\rho(1))
\] is unitary, even if \(\rho\) is degenerate, and \([P_F,U]\) is
compact. Thus \(P_FUP_F\) on \(P_F(H\oplus H)\) is Fredholm:
\(P_FU^*P_F\) is an inverse modulo compacts, since the errors factor
through \([P_F,U]\). Define \(I_u\) to be its index.

Unitary equivalence preserves it. The normalization respects direct sums
up to a permutation of the doubled summands, so its index is additive.
For a degenerate cycle it gives zero. To see this last point without an
extra convention, let \(H_e=\overline{\pi(A)H}\). Exact commutation
makes \(H_e\) reducing for \(F\); the zero self-adjointness and square
defects make \(F|_{H_e}\) a self-adjoint involution. The normalization
leaves that restriction unchanged and has zero off-diagonal \(s\) there,
while \(\pi\) is zero on \(H_e^\perp\). Consequently \(N(F)\) commutes
exactly with \(\rho(A)\), so \(P_FUP_F\) is unitary on its range and has
index zero.

Under an operator homotopy \(F_t\), the normalized projections \(P_t\)
vary in norm. Their ranges can be locally identified by norm-continuous
unitaries. Here is the formula. If \(\|P-Q\|<1\), set \[
X=QP+(1-Q)(1-P),\qquad V=X(X^*X)^{-1/2}.
\] The identities \(XP=QX\) and \(X^*X=1-(P-Q)^2\) show that \(V\) is a
unitary with \(VPV^*=Q\); it depends continuously on \(Q\). Cover a
compact parameter interval by finitely many such neighborhoods and
concatenate the identifications. After transporting the ranges by these
unitaries, the compressed operators are a norm-continuous Fredholm path
on a fixed space. Its index is constant. This proves invariance under
every generating relation, hence well-definedness and additivity on
\(K^1(A)\).

If the starting \(F\) is already an exact involution, its normalization
is \(\operatorname{diag}(F,-F)\). The second copy has zero
representation and contributes the identity compression, of index zero.
Removing a zero-representation part of a unital algebra action likewise
contributes zero. Thus the displayed formula agrees with the
construction. \(\square\)

Applying the lemma to the Hardy module and \(u=z\), (5.2) gives
\(I_z=-1\). In particular its class has infinite order. This detects the
class directly in the equivalence relation of this lesson, without
assuming the later class-level extension isomorphism.

\subsection{8. Exercises}\label{exercises}

\textbf{Exercise 8.1 (basic).} Verify that a self-adjoint unitary
commuting with \(\pi(A)\) is degenerate. In the even case include the
grading condition. Explain why \(F=0\) on a finite-dimensional nonzero
representation is usually not degenerate.

\textbf{Exercise 8.2 (intermediate).} Give the operator homotopy between
two cycle operators differing by a compact operator. Check its square
and adjoint defects explicitly.

\textbf{Exercise 8.3 (intermediate).} Verify all the Hardy-module
conditions by computing \([P,\pi(z^k)]\) on Fourier basis vectors.
Identify the compression for \(z^{-1}\).

\textbf{Exercise 8.4 (advanced).} Prove that the Hardy class has
infinite order, using the positive-compression index. Determine the
class detected after replacing \(F\) by \(-F\).

\subsection{9. Solutions}\label{solutions}

\textbf{Solution 8.1.} Self-adjointness makes \(F-F^*=0\), involutivity
makes \(F^2-1=0\), and commutation makes \([F,\pi(a)]=0\). Thus all
three defects vanish. An even cycle must also have \(\Gamma F=-F\Gamma\)
and an even representation; the conclusion applies once those cycle
requirements hold. For a finite-dimensional example with \(F=0\), the
square defect is \(-\pi(a)\). It is compact because the space is finite
dimensional, but it is not zero when \(\pi(a)\ne0\). Being a cycle and
being degenerate are distinct conditions.

\textbf{Solution 8.2.} Put \(K=F'-F\in\mathcal K(H)\) and \(F_t=F+tK\).
Then \[
\begin{aligned}
\left[F_t,\pi(a)\right]&=[F,\pi(a)]+t[K,\pi(a)],\\
F_t-F_t^*&=F-F^*+t(K-K^*),\\
F_t^2-1&=F^2-1+t(FK+KF)+t^2K^2.
\end{aligned}
\] Every added term is compact, so multiplying the last two expressions
by \(\pi(a)\) preserves compactness. The path is norm-continuous. If
both endpoints are odd for the grading, then \(K\) and every \(F_t\) are
odd. This verifies an operator homotopy directly. The same conclusion
for a locally compact difference follows from the ideal calculation of
Proposition 2.2.

\textbf{Solution 8.3.} For \(k\geq0\), \[
[P,\pi(z^k)]e_n
=\big(1_{\{n+k\geq0\}}-1_{\{n\geq0\}}\big)e_{n+k}.
\] It is nonzero exactly for \(-k\leq n\leq-1\), and then equals
\(e_{n+k}\). For \(k=-\ell<0\), it is nonzero exactly for
\(0\leq n\leq\ell-1\), and equals \(-e_{n-\ell}\). Thus its rank is
\(|k|\), with rank zero when \(k=0\). Linear combinations give
finite-rank commutators for trigonometric polynomials, and uniform
approximation gives compact commutators for all continuous functions.
The other defects vanish because \(F=2P-1\) is self-adjoint and squares
to one. On the nonnegative basis, compression of \(\pi(z^{-1})\) is
\(S^*\): it kills \(e_0\) and sends \(e_n\) to \(e_{n-1}\) for
\(n\geq1\).

\textbf{Solution 8.4.} Lemma 7.1 gives a homomorphism \(I_z\) on
\(K^1(C(S^1))\). For the Hardy class \(h\), its value is \(-1\), so
\(I_z(nh)=-n\). If \(nh=0\), this forces \(n=0\); hence \(h\) has
infinite order. The replacement \(-F\) represents its inverse by Theorem
3.1. Directly, its positive projection is \(1-P\). Compression of
multiplication by \(z\) to the negative Fourier modes kills \(e_{-1}\),
sends \(e_{-j}\) to \(e_{-(j-1)}\) for \(j\geq2\), and is surjective on
that negative subspace. Its index is \(+1\), agreeing with the inverse
class and fixing the sign.

\subsection{Internal proof dependencies and source
credit}\label{internal-proof-dependencies-and-source-credit}

The Hilbert-space Fredholm facts, including composition additivity, are
proved in Lemmas 1.1--1.2 of the extension lesson. Continuous functional
calculus is proved in
\href{https://kokunoyumeto.github.io/open-math-courses/courses/foundations-of-von-neumann-algebras/c-star-algebras-continuous-functional-calculus-automatic-continuity-positive-cones.html}{C*-algebras:
continuous functional calculus, automatic continuity, positive cones,
approximate identities and quotients}, Theorems 5.1 and 6.1. Lemma 4.0
proves the compact-resolvent spectral domain and calculus used in
Theorem 4.1 here, without taking an unbounded spectral theorem as an
external input. The finite-chart Sobolev construction and compactness,
frozen-coefficient elliptic estimate, and uniformly bounded, strongly
vanishing Friedrichs commutator are proved here in Lemmas 4.1a--4.1c.
Corollary 4.2 uses these proofs to identify the distributional adjoint
domain with \(H^1\) and prove self-adjointness and compact resolvent for
every formally self-adjoint first-order elliptic operator on a closed
manifold. The exact earlier measure and Hilbert-space inputs are linked
before Lemma 4.1a; no pseudodifferential parametrix or external
elliptic-regularity theorem supplies this bridge.

The class-level isomorphisms (6.3) and their precise normalization,
dilation and homotopy assertions remain proof obligations in
\emph{Pictures of KK: Fredholm picture, quasihomomorphisms, \(KK^1\) and
extensions}. The localized noncompact bounded-transform theorem remains
a proof obligation in \emph{Unbounded Kasparov modules and spectral
triples}. The cycle constructions, group laws, compact-resolvent bounded
transform and circle detection are proved above from those existing
internal prerequisites. The full class-level identifications in (6.3)
and the localized noncompact extension remain unproved here.

\subsection{References}\label{references}

\begin{itemize}
\item
  B. Blackadar, \emph{K-Theory for Operator Algebras}, freely accessible
  author-corrected second-edition PDF. Sections 17.1--17.5 supply the
  cycle, normalization, group and Fredholm pictures; Section 17.11
  treats the bounded transform.
  \href{https://www.bruceblackadar.com/Mathematics/book6.pdf}{Free
  author version}. The resolvent and compact spectral proofs required
  here are written in Lemma 4.0 and Theorem 4.1.
\item
  B. Blackadar, \emph{Operator Algebras: Theory of C*-Algebras and von
  Neumann Algebras}, freely accessible author-revised version dated 8
  February 2017, I.6.1 and I.8.4.
  \href{https://www.bruceblackadar.com/Mathematics/Cycr.pdf}{Free
  revised author version}. Lemma 4.0 supplies the compact positive
  spectral argument rather than citing it in place of proof.
\item
  J. Ruppenthal, \emph{Friedrichs' Extension Lemma with Boundary Values
  and Applications in Complex Analysis}, free author preprint dated 14
  October 2009, Theorem 3.1 and Lemma 3.2, pp.~5--6.
  \href{https://www2.math.uni-wuppertal.de/~ruppenth/papers/friedrichs14102009.pdf}{Actual
  free author PDF}. The matrix-kernel bound, strong convergence and
  weak-domain argument required here are fully proved in Lemma 4.1c.
\item
  C. Bär and W. Ballmann, \emph{Boundary Value Problems for Elliptic
  Differential Operators of First Order}, arXiv:1101.1196v2, Sections
  2.1--2.2 and Lemma 3.1/Remark 3.2.
  \href{https://arxiv.org/abs/1101.1196}{Actual free preprint}. These
  sections provide the source comparison for measured manifolds, formal
  adjoints and maximal domains. The Sobolev, regularity and compactness
  proofs used here are local Lemmas 4.1a--4.1c.
\end{itemize}

These free source PDFs retain their authors' copyrights. The CC0 notice
applies to this lesson's original text.

\end{document}
